\documentclass[11pt]{article}
\usepackage[T1]{fontenc}
\usepackage{lmodern,amsmath,amssymb,amsthm,mathtools,booktabs,microtype}
\usepackage[margin=1in]{geometry}
\usepackage[colorlinks=true,linkcolor=blue!55!black,urlcolor=blue!55!black,citecolor=blue!55!black]{hyperref}
\usepackage{xcolor}
\newtheorem{theorem}{Theorem}[section]
\newtheorem{lemma}[theorem]{Lemma}
\newtheorem{corollary}[theorem]{Corollary}
\newtheorem{proposition}[theorem]{Proposition}
\theoremstyle{remark}\newtheorem{remark}[theorem]{Remark}
\newcommand{\E}{\mathbb E}
\newcommand{\R}{\mathbb R}
\newcommand{\Z}{\mathbb Z}
\newcommand{\U}{\mathcal U}
\newcommand{\bias}{\operatorname{bias}}
\newcommand{\supp}{\operatorname{supp}}
\newcommand{\bigO}{\mathcal O}
\newcommand{\D}{\mathcal D}
\newcommand{\B}{\mathcal B}
\title{The missing Hermite constant:\\ sharp small-bias thresholds, lattice oscillations,\\ and an application to quadratic prime patterns}
\author{Research note for the ProveIt project}
\date{27 November 2026}
\begin{document}
\maketitle
\begin{abstract}
The merged ProveIt article asks for the subsequential leading constants of the minimum parity bias of a law on the Boolean cube that misses a Hamming cap at a Hermite threshold. We determine all of them. If $h_j$ is the largest zero of the monic probabilists' Hermite polynomial $H_j$, and $\delta_{k,j}$ is the distance from $h_j\sqrt{k}$ down to the nearest integer of parity $k$, then the minimum bias is
\[
 b_k(h_j)=j!k^{-j}(1+\delta_{k,j}+O_j(k^{-1/2})).
\]
For $j\ge2$ the normalized values are asymptotically uniformly distributed on $[j!,3j!]$. We prove a uniform left critical-window formula and a rigidity theorem for near optimizers, without any Riemann hypothesis. A parity-sensitive separating polynomial and a conductor-weighted estimate transfer the result to primes realizing prescribed quadratic signs, using only Dirichlet GRH as stated in the vendored Lean development. Exact rational finite-dimensional lower and upper certificates accompany the paper. The new proofs are mathematical proofs, not yet Lean formalizations.
\end{abstract}

\section{Source, scope, and contribution}
The source is the vendored readable edition of \texttt{lsernau/grh-lean} in ProveIt \cite{vendor,upstream}, whose README pins upstream commit
\begin{center}\small\texttt{72fead102f663d5c7a407964a13dc0cf9d977207}.\end{center}
We inspected its README and the declarations in \texttt{src/Statements.lean} and \texttt{src/Main.lean}. In particular, \texttt{LFunction\_ne\_zero\_of\_half\_lt\_re} asserts nonvanishing for Dirichlet $L$-functions in $\Re s>1/2$, except for the principal pole, and \texttt{LFunction\_eq\_zero\_re} locates all open-strip zeros on the critical line. The upstream README reports builds and independent proof checks with the three standard Lean axioms. We use these statements as the input theorem (Dirichlet GRH); we have not rerun those checks. No nonabelian GRH is assumed or concluded.

The article \cite{article} (131 reports, 723 pages, 914 theorem-like statements according to its README at retrieval) distinguishes results that are unconditional, use Dirichlet GRH, or reduce a problem to an open estimate. Its source section describes the passage from one-sided bounds on smoothed prime sums to zero exclusion. Here we use the theorem in the opposite direction, via the smoothed prime explicit formula, and extend the article's Fourier extremal approach in \texttt{sections/dir.tex}. Specifically, the lemma \texttt{dir:lem:hamming}(b) proves $b_k(h_j)\asymp_j k^{-j}$; \texttt{dir:prop:hammingdual} identifies the dual Littlewood complexity; \texttt{dir:thm:signs} gives the prime application. Open Question 14 asks for the leading constants at $h_2,h_3,\ldots$ and the critical window. We answer the first in full and the second on the left, uniformly up to the threshold itself. The right-hand window is not claimed. The new ingredient is the top two layers' different parity, both in the separating polynomial and in the extremizing measure. Related literature includes Bella\"iche's Littlewood complexity \cite{bellaiche} and the work of Derksen--Ivanov--Lee--Viola on symmetric small-bias distributions \cite{dilv}. A limited literature search did not establish prior occurrence of our formula; global publication priority is not asserted.

\section{Definitions and exact asymptotic}
Let $k\ge1$, $u=(u_1,\ldots,u_k)\in\{-1,1\}^k$, $S=\sum_i u_i$, $X=S/\sqrt{k}$, and $w=(k-S)/2$. For $A\subseteq[k]$, write $u_A=\prod_{i\in A}u_i$, with $u_\varnothing=1$. For a probability law $\mu$ on the cube set
\[
 \bias(\mu)=\max_{A\ne\varnothing}|\E_\mu u_A|,\quad
 b_k(a)=\min_{\mu:\,\supp\mu\subseteq\{X\le a\}}\bias(\mu).
\]
The minimum exists by compactness whenever the allowed set is nonempty. For a real function $g$ use Walsh coefficients $\widehat g(A)=\E_\U g u_A$, where $\U$ is uniform. Write $L(g)=\sum_A|\widehat g(A)|$ and $L_*(g)=L(g)-|\widehat g(\varnothing)|$. The Fourier expansion implies
\begin{equation}\label{eq:dualbound}
 |\E_\mu g-\E_\U g|\le \bias(\mu)L_*(g),\qquad L(gh)\le L(g)L(h).
\end{equation}
Let $H_0=1,H_1=x,H_{j+1}=xH_j-jH_{j-1}$, and $h_j$ be the largest zero of $H_j$. For $j=1,2,3,4$ the values are $0,1,\sqrt3,\sqrt{3+\sqrt6}$.

\begin{theorem}[Sharp lattice threshold]\label{thm:main}
Fix $j\ge1$ and $R>0$. Let $s\equiv k\pmod2$ be an integer such that
\[
 r=s-h_j\sqrt{k}\in[-R,0].
\]
Uniformly in such $s$, as $k\to\infty$,
\begin{equation}\label{eq:main}
 b_k(s/\sqrt{k})=k^{-j}\bigl(j!(1-r)+O_{j,R}(k^{-1/2})\bigr).
\end{equation}
The upper bound is attained to this accuracy by a permutation-invariant law on at most $4j-2$ Hamming layers. For $\lambda\le0$ fixed let $a_k=h_j+\lambda/\sqrt{k}$ and $\delta_k= h_j\sqrt{k}+\lambda-s_k$ with $s_k$ the largest integer of parity $k$ at most $h_j\sqrt{k}+\lambda$. Then
\begin{equation}\label{eq:window}
 k^j b_k(a_k)=j!(1-\lambda+\delta_k)+O_{j,\lambda}(k^{-1/2}),\qquad 0\le\delta_k<2.
\end{equation}
The error is uniform for $\lambda$ in any fixed compact subinterval of $(-\infty,0]$.
\end{theorem}

The proof will also give the structure of near minimizers. The fact that the minimum is over \emph{all} laws, not just symmetric ones, is essential for the prime application.

\section{Hermite and layer preliminaries}
Throughout put $h=h_j$, $J=j!$, and $Q(x)=H_j(x)/(x-h)$, a monic polynomial of degree $j-1$. For $G$ a standard normal variable, Hermite orthogonality gives $\E H_j(G)p(G)=0$ if $\deg p<j$ and $\E H_j(G)^2=J$. These follow by integration by parts from Rodrigues' formula. We have
\begin{equation}\label{eq:hermite}
 F(x)=(x-h)Q(x)^2=H_j(x)Q(x),\quad \E F(G)=0,\quad \E Q(G)^2=J.
\end{equation}
For the last identity, use $H_j'=jH_{j-1}=Q+(x-h)Q'$. Thus $\E H_j'(G)Q(G)=j\,(j-1)!=J$ by monicity and orthogonality, while $\E H_j(G)Q'(G)=0$. For any fixed polynomial $p$ the cube moments satisfy
\begin{equation}\label{eq:moment}
 \E_\U p(X)=\E p(G)+O_p(k^{-1}).
\end{equation}
Indeed odd moments vanish, and for degree $2t$ the tuples of $t$ distinct pairs contribute $(2t-1)!!(k)_t/k^t$, with all remaining even-index tuples contributing $O_t(k^{-1})$. Here $(k)_t=k(k-1)\cdots(k-t+1)$.

\begin{lemma}[Quadrature with an exact top node]\label{lem:quad}
For $a$ in a fixed sufficiently small neighborhood of $h$, there is a probability law $\nu_a$ with $j$ simple real nodes, largest node $a$, and positive masses, whose moments through degree $2j-2$ equal those of $G$. As $a\to h$, the nodes and weights differ from those of Gaussian quadrature $\nu_h$ by $O_j(|a-h|)$. In particular, if $w_*$ is the top-node mass of $\nu_h$, then $w_*Q(h)^2=J$.
\end{lemma}
\begin{proof}
Consider the $j\times j$ real symmetric Jacobi matrix for multiplication by $x$ in the basis $H_\ell/\sqrt{\ell!}$, $0\le\ell<j$, with off-diagonal entries $\sqrt{1},\ldots,\sqrt{j-1}$. Modify its last diagonal entry to $c=H_j(a)/H_{j-1}(a)$. The determinant recurrence gives characteristic polynomial $H_j-cH_{j-1}$; the denominator is nonzero near $h$ by strict interlacing (or irreducibility of the tridiagonal matrix). For $a$ near $h$ this prescribed eigenvalue is the largest, since all eigenvalues at $c=0$ are simple and separated. Simplicity follows since an eigenvector is determined by its first entry, which cannot be zero. The spectral law at the first coordinate has positive weights. The walk expansion of the $\ell$th moment agrees with that of the infinite Hermite Jacobi matrix through $\ell=2j-2$: the last diagonal entry cannot be used and returned from in fewer than $2j-1$ steps, and the first omitted vertex cannot be visited and returned from in fewer than $2j$ steps. For $j=1$ this just says the mass is at $a$ and matches the constant moment. The simple eigenvalues and their spectral projections are smooth in $c$ near zero (e.g. implicit differentiation of the characteristic polynomial and the spectral projection product formula), and $c=O_j(|a-h|)$. Finally $Q$ vanishes at every node of $\nu_h$ except $h$, and the moment of $Q^2$ is exact, so \eqref{eq:hermite} yields the mass identity.
\end{proof}

For a layer $w$, define $R_{\ell,k}(w)$ as the mean of $u_A$ on the uniform layer for any $|A|=\ell$. Then
\begin{align}
 \sum_{\ell=0}^k\binom{k}{\ell} R_{\ell,k}(w)z^\ell&=(1-z)^w(1+z)^{k-w},\
 R_{k-\ell,k}(w)&=(-1)^w R_{\ell,k}(w).\label{eq:reflect}
\end{align}
With $P_{\ell,k}(X)=(k)_\ell k^{-\ell/2}R_{\ell,k}(w)$, multiplication by $S$ gives the monic polynomial recurrence
\[
 P_{0,k}=1,\quad P_{1,k}=X,\quad
 P_{\ell+1,k}=XP_{\ell,k}-\ell(1-(\ell-1)/k)P_{\ell-1,k}.
\]
Consequently $P_{\ell,k}=H_\ell+O_\ell(k^{-1})$ coefficientwise for fixed $\ell$, and $\E_\U P_{\ell,k}=0$ for $\ell\ge1$. Also the normalized layer sum of all Walsh characters of degree $\ell$ has squared uniform norm $1/\binom{k}{\ell}$ by orthogonality. At any layer with $|X|\le C$ the uniform mass is $\gg_C k^{-1/2}$ by Stirling's formula. Hence
\begin{equation}\label{eq:middle}
 |R_{\ell,k}(w)|\ll_C k^{1/4}\binom{k}{\ell}^{-1/2}\quad (|X|\le C).
\end{equation}

\section{Proof of the sharp threshold}
\subsection{Lower bound: a parity-sensitive polynomial}
Set $w_0=(k-s)/2$ and $\eta(u)=(-1)^{w-w_0}=(-1)^{w_0}\prod_i u_i$. If $S\le s$ then
\begin{equation}\label{eq:parity}
 S\le s-1+\eta(u).
\end{equation}
In fact, when $\eta=1$ the maximum is $s$, and when $\eta=-1$ it is $s-2$. This is the lattice effect missed by a single continuous cutoff. Define the witness
\begin{equation}\label{eq:witness}
 g(u)=\left(X-h+\frac{1-r-\eta(u)}{\sqrt{k}}\right)Q(X)^2.
\end{equation}
By \eqref{eq:parity}, $g\le0$ on the entire permitted set, independently of the law. For $k>2j-2$, the uniform mean of $\eta Q(X)^2$ vanishes because $Q(X)^2$ has Walsh degree at most $2j-2$. Thus \eqref{eq:hermite}--\eqref{eq:moment} give
\begin{equation}\label{eq:gmean}
 \E_\U g=\frac{J(1-r)}{\sqrt{k}}+O_{j,R}(k^{-1}).
\end{equation}
Since $L(X)=\sqrt{k}$ and $L(\eta)=1$, monicity of $F$ and the fixed coefficients imply
\begin{equation}\label{eq:gcost}
 L_*(g)\le L(g)\le k^{j-1/2}+O_{j,R}(k^{j-1}).
\end{equation}
For every law supported on $S\le s$, \eqref{eq:dualbound} gives $\bias(\mu)\ge \E_\U g/L_*(g)$. The numerator is positive for large $k$. Dividing \eqref{eq:gmean} by \eqref{eq:gcost} proves the lower bound in \eqref{eq:main} with error $O_{j,R}(k^{-j-1/2})$.

\subsection{Upper bound: keep the two top nodes exact}
Split the permitted layers into the parity classes $\eta=1$ and $\eta=-1$, with respective largest nodes $a_+=(s/\sqrt{k})=h+r/\sqrt{k}$ and $a_-=(s-2)/\sqrt{k}=h+(r-2)/\sqrt{k}$. Put $r_+=r$ and $r_-=r-2$. The grid spacing in each class is $4/\sqrt{k}$.

For each class use the law of Lemma~\ref{lem:quad} with largest node $a_\pm$. Leave this node unchanged. For every other node choose one grid point on each side at distances in $[4/\sqrt{k},8/\sqrt{k}]$ (if an endpoint is a grid point either admissible choice suffices). For large $k$ the points are distinct, in the class, and below its largest permitted node. Split the original node mass between its two grid points by linear interpolation. The weights are bounded below by a positive constant depending on $j$, since the splitting fractions are at least $1/3$ and the quadrature weights converge to positive limits. The moments through $2j-2$ now differ from those of the Gaussian by $O_{j,R}(k^{-1})$, by Taylor expansion; the uniform cube moments have the same accuracy by \eqref{eq:moment}.

There are exactly $2j-1$ grid nodes per class. Correct their weights to match the cube moments of degrees $0,\ldots,2j-2$ exactly. The square Vandermonde system is invertible. Its inverse entries are $O_{j,R}(\sqrt{k})$: each row consists of coefficients of a Lagrange basis polynomial; each denominator has at most one separation of order $k^{-1/2}$ (the nearby paired point), and all other separations are bounded away from zero. Hence the weight correction is $O_{j,R}(k^{-1/2})$, retaining positivity and total mass one. For $j=1$ the system is just the constant moment, and the top node is the entire law. Denote by $\mu_\pm$ the resulting distributions obtained by spreading the weights uniformly on their layers, and take $\mu=(\mu_++\mu_-)/2$.

Every low degree $1\le\ell\le 2j-2$ has exactly zero mean under both classes by moment matching. At degree $2j-1$, $P_{2j-1,k}-F$ has degree at most $2j-2$, so
\begin{equation}\label{eq:odd}
 \E_{\mu_\pm} P_{2j-1,k}=\E_{\mu_\pm} F-\E_\U F
 =\frac{Jr_\pm}{\sqrt{k}}+O_{j,R}(k^{-1}).
\end{equation}
Indeed at every interior node $F$ has a double zero in the limiting quadrature, so it is $O(k^{-1})$ on the paired points. At the top node its contribution is $(r_\pm/\sqrt{k})(w_* Q(h)^2+O(k^{-1/2}))$. Lemma~\ref{lem:quad} evaluates the coefficient. Similarly,
\begin{equation}\label{eq:even}
 \E_{\mu_\pm}P_{2j,k}=-J+O_{j,R}(k^{-1/2}).
\end{equation}
To see this, the limiting law is $\nu_h$, and $P_{2j,k}$ tends to $H_{2j}$. The polynomial $H_{2j}-H_j^2$ has degree at most $2j-2$ (both polynomials are monic and even), so its quadrature mean equals its Gaussian mean $-J$. Also $H_j^2$ vanishes on the limiting quadrature nodes. The convergence rate follows from the node and weight bounds already established.

Using $(k)_\ell=k^\ell(1+O_\ell(k^{-1}))$ and the reflection identity \eqref{eq:reflect}, the four potentially largest nontrivial coefficients of the mixture satisfy
\begin{center}
\begin{tabular}{@{}cl@{}}
\toprule
Degree & $k^j$ times the common coefficient at that degree\\\midrule
$2j-1$ & $J(r-1)+O(k^{-1/2})$\\
$k-(2j-1)$ & $(-1)^{w_0}J+O(k^{-1/2})$\\
$2j$ & $-J+O(k^{-1/2})$\\
$k-2j$ & $O(k^{-1/2})$\\\bottomrule
\end{tabular}
\end{center}
The top degree $k$ has mean zero by equal parity masses; low degrees and their complements already vanish. The remaining degrees $2j+1\le\ell\le k-2j-1$ have absolute bias $O_j(k^{-j-1/4})$ by \eqref{eq:middle}, since all nodes lie in a fixed compact interval and $\binom{k}{\ell}\ge\binom{k}{2j+1}$. This is strictly smaller than the leading bound for sufficiently large $k$. Because $r\le0$, $J(1-r)\ge J$, proving the matching upper bound, uniformly over the specified range. Substitution $r=\lambda-\delta_k$ establishes the critical-window assertion. \qed

\section{Consequences: oscillation, duality, rigidity}
\begin{corollary}[All subsequential constants]\label{cor:limits}
For $j\ge2$ and fixed $\lambda\le0$, the empirical distribution of $k^j b_k(h_j+\lambda/\sqrt{k})$ for $1\le k\le N$ converges, as $N\to\infty$, to the uniform law on
\[
 [j!(1-\lambda),\ j!(3-\lambda)].
\]
In particular, for $\lambda=0$ its liminf and limsup are $j!$ and $3j!$, and every number between is a subsequential limit. For $j=1,\lambda=0$, the two parity subsequences have limits $1$ and $2$ respectively (even and odd $k$).
\end{corollary}
\begin{proof}
Fix the parity $e=0$ or $1$ and write $k=2n+e$. Then $\delta_k=2\{(h_j\sqrt{2n+e}+\lambda-e)/2\}$, with braces denoting fractional part. For $j\ge2$, $h_j>0$. The fractional parts of $a\sqrt{2n+e}+c$ are uniformly distributed for any fixed $a>0,c\in\R$. For completeness, for an interval $[v,w)\subseteq[0,1)$, invert the inequalities $m+v\le a\sqrt{2n+e}+c<m+w$ over each complete integer block $m$. The number of integers $n$ in a block differs from its interval length by at most two; the interval length equals $(w-v)$ times the full block length plus $O_{a,c}(1)$, by expansion of the quadratic inverse. There are $O(\sqrt{N})$ blocks and the incomplete blocks have total length $O(\sqrt{N})$, so the empirical proportion equals $w-v+O_{a,c}(N^{-1/2})$. Combining parities and \eqref{eq:window} yields weak convergence to the claimed law; each endpoint is approximated infinitely often by the same counting argument. For $j=1$, $h_1=0$ and $\delta_k=0$ or $1$ according to parity, as stated.
\end{proof}

For $j=2$ the cluster interval is $[2,6]$; for $j=3$ it is $[6,18]$; for $j=4$ it is $[24,72]$. In particular there is no universal single leading constant at these thresholds. Along $k=n^2$ for $j=2$, $s=n$ is of the correct parity, so the lower endpoint $2$ occurs; along $k=n^2+2$ the parity is still that of $n$ and the residual tends to zero, while along $k=n^2-2$ the top permitted integer is $n-2$ and the residual tends to two, producing the upper endpoint $6$.

\begin{corollary}[Optimized Littlewood complexity]\label{cor:littlewood}
Let $D=\{S>s\}$ and define $\Phi(D)$ as the infimum of $L(f)/\E_\U f$ over real functions $f$ with $\E_\U f>0$ and $f\le0$ off $D$. Under the hypotheses of Theorem~\ref{thm:main},
\[
 \Phi(D)=1+\frac1{b_k(s/\sqrt{k})}
 =\frac{k^j}{j!(1-r)}(1+O_{j,R}(k^{-1/2}))+1.
\]
\end{corollary}
\begin{proof}
The vector $v(u)=(u_A)_{A\ne\varnothing}$ has finite dimension. The least bias is the $\ell^\infty$ distance of zero from the convex hull of the $v(u)$ with $u\notin D$. Duality of $\ell^\infty$ and $\ell^1$ and finite minimax identify it with $\max_{\|z\|_1\le1}\min_{u\notin D}z\cdot v(u)$. It is positive, since a law with all nonconstant Walsh coefficients zero is uniform, which is not supported off a nonempty $D$. An optimal $z$ gives the witness $b_k-z\cdot v$, and conversely \eqref{eq:dualbound} bounds any witness's nonconstant Fourier norm below by its mean divided by $b_k$. This proves equality; \eqref{eq:main} supplies the expansion.
\end{proof}

\begin{theorem}[Top-layer rigidity]\label{thm:rigidity}
In the setting of Theorem~\ref{thm:main}, suppose a sequence of arbitrary laws $\mu_k$ supported on $S\le s$ satisfies
\[
 k^j\bias(\mu_k)\le J(1-r)+o(1).
\]
Then $X$ converges in distribution under $\mu_k$ to $\nu_h$. More precisely, each parity class has mass $1/2+o(1)$, its conditional distribution of $X$ tends to $\nu_h$, and
\[
 \mu_k(S=s)\longrightarrow w_*/2,\qquad
 \mu_k(S=s-2)\longrightarrow w_*/2.
\]
No permutation symmetry is required.
\end{theorem}
\begin{proof}
The witness \eqref{eq:witness} is nonpositive on the support. Equations \eqref{eq:dualbound}, \eqref{eq:gmean} and \eqref{eq:gcost} imply
\[
 0\le\E_{\mu_k}[-g]\le \bias(\mu_k)L_*(g)-\E_\U g=o(k^{-1/2}).
\]
For every fixed degree $t\le 2j$, $L(X^t)\le k^{t/2}$, so the moments of $X$ through degree $2j-1$ tend to the Gaussian moments and the even moment $2j$ stays bounded. The same statements hold conditional on either parity: its indicator is $(1\pm\eta)/2$, its mass tends to $1/2$ by the bias bound on the top character, and $L(\eta X^t)\le k^{t/2}$ with zero uniform mean for $t<k$. In particular the conditional laws are tight, and moments through degree $2j-1$ are uniformly integrable. On every compact set, $g$ differs from $F$ uniformly by $O(k^{-1/2})$. Any weak subsequential limit is supported on $(-\infty,h]$, and the nonnegative function $-F$ there is positive except at the $j$ zeros of $H_j$. The vanishing expectation of $-g$ forces the limit to be supported on those zeros (apply the bound to compact neighborhoods away from the zeros). Its moments through degree $j-1$ determine the weights uniquely by a Vandermonde system and are Gaussian. Thus it equals $\nu_h$ for each parity. Finally
\[
 \sqrt{k}(-g)=(s-1+\eta-S)Q(X)^2.
\]
On a fixed small neighborhood of $h$, $Q(X)^2$ is bounded below by a positive constant, and if $S$ is not the top layer of its parity the first factor is at least four. Its expectation tends to zero. The neighborhood has limiting conditional mass $w_*$; all but $o(1)$ of this mass lies on the exact top layer. Multiplication by the limiting parity mass proves both assertions.
\end{proof}

\section{Transfer to primes using Dirichlet GRH}
Let $q_1,\ldots,q_k\le Y$ be distinct odd primes and $v\in\{-1,1\}^k$. For a prime $p$ not dividing $\mathcal Q=\prod_i q_i$, let $u_i(p)=v_i(p/q_i)$ with Legendre symbol, and $M_v(p)=(k+S(u(p)))/2$. We use the standard once-smoothed prime explicit formula under Dirichlet GRH, in the form used in \cite[\texttt{dir:thm:signs}]{article} (see also \cite{bellaiche,mv}). For $A\subseteq[k]$, $f_A=\prod_{i\in A}q_i$ and $\chi_A=\prod_{i\in A}(\cdot/q_i)$, primitive of conductor $f_A$ by the Chinese remainder theorem. For $x\ge3$,
\begin{equation}\label{eq:prime}
 \sum_{p<x,\ p\nmid\mathcal Q}(1-p/x)\log p\,\chi_A(p)
 =\mathbf1_{A=\varnothing}\frac{x}{2}
 +O\bigl(\sqrt{x}\log(2f_A)+\log\mathcal Q\bigr).
\end{equation}
The absolute implied constant is uniform. The kernel $x^z/(z(z+1))$ makes the zero sum absolutely convergent; the zero-counting bound is $O(T\log(f_A(T+2)))$ and GRH bounds the numerator by $\sqrt{x}$. The pole contributes $x/2$ for the principal character; the trivial-zero and contour terms are included in the displayed error. Removing prime powers costs $O(\sqrt{x})$ using the elementary Chebyshev estimate, and deleting primes dividing $\mathcal Q$ costs at most $\log\mathcal Q$. Thus only Dirichlet GRH, not a general Artin or Hecke hypothesis, enters.

\begin{theorem}[Critical-window prime patterns]\label{thm:prime}
Assume the source Dirichlet GRH theorem. For fixed $j\ge1$ and $R>0$, and every sufficiently large $k$, take $s$ as in Theorem~\ref{thm:main}. There is a prime $p\nmid\mathcal Q$ satisfying
\[
 M_v(p)>\frac{k+s}{2},\qquad
 p\ll_{j,R}\frac{k^{2j}\log^2(2Y)}{(1-r)^2}.
\]
In particular this is uniform for the left critical window $a_k=h_j+\lambda/\sqrt{k}$, and $M_v(p)>(k+h_j\sqrt{k}+\lambda)/2$. The exponent $2j$ is also the exponent predicted by the sharp optimized unweighted Littlewood complexity; the theorem does not assert an arithmetic lower bound for the least prime.
\end{theorem}
\begin{proof}
Use the explicit witness $g$ in \eqref{eq:witness}. Its coefficients split into a low-degree part of degree at most $2j-1$ with Fourier norm $O_{j,R}(k^{j-1/2})$, and the high-degree part $-\eta Q(X)^2/\sqrt{k}$ of norm $O_j(k^{j-3/2})$. Consequently, despite the latter involving characters of conductors as large as $\mathcal Q$,
\[
 \sum_A |\widehat g(A)|\log(2f_A)\ll_{j,R} k^{j-1/2}\log(2Y).
\]
For the low part $\log(2f_A)\ll_j\log(2Y)$; for the high part it is at most $k\log(2Y)$. The translated coefficients for the signs $v_i$ have the same absolute values. Multiply \eqref{eq:prime} by those coefficients and sum:
\[
 \sum_{p<x,\ p\nmid\mathcal Q}(1-p/x)\log p\,g(u(p))
 =\frac{x}{2}\E_\U g+O_{j,R}\left(k^{j-1/2}(\sqrt{x}+k)\log(2Y)\right).
\]
By \eqref{eq:gmean}, the mean is at least $J(1-r)/(2\sqrt{k})$ for large $k$. Choose $x=C_{j,R} k^{2j}\log^2(2Y)/(1-r)^2$. The main term dominates both error terms for a sufficiently large constant $C_{j,R}$ (use $j\ge1$, $1\le1-r\le1+R$ and $\log(2Y)\ge1$). The sum is then positive, so some prime has $g(u(p))>0$, which by \eqref{eq:parity} implies $S>s$. Since $S$ has parity $k$, exceeding the rounded threshold is equivalent to exceeding the original real threshold in the critical window.
\end{proof}

The exponent at an exact root was already proved in the article. The extension here is the uniform critical-window statement with a parity-sensitive, asymptotically optimal unweighted separator; its expensive high-degree part has just the right suppressed Fourier mass to preserve the prime bound. The new arithmetic statement is not a claim of an improved exponent or of an explicit best leading prime constant.

\section{Verification, limitations, and further questions}
The accompanying Python program computes exact rational lower and upper certificates for $j=2,3$. It does not solve the linear program numerically: a rational witness is checked by its Walsh coefficients, and rational layer weights are checked by moments and every nontrivial parity coefficient using integer Krawtchouk recurrences. The output bounds the finite-dimensional optimum from both sides. It is supporting evidence, not a replacement for the proof. See \texttt{README.md}, \texttt{check.py}, and \texttt{checks.txt}. The proofs in this note are self-contained apart from standard Gaussian/Hermite and smoothed explicit-formula inputs; they have not been checked in Lean. The source files were inspected using the GitHub connector, not the entire 818-file Lean dependency graph. The article's other sections were surveyed for overlap, not independently validated. No repository files were changed.

Further directions: (i) determine the right critical-window function when $r>0$, where optimizing fourth or higher even moments may allow rare extreme layers; (ii) optimize the conductor-weighted Fourier norm instead of the unweighted norm, to obtain sharp prime constants or identify a barrier to improvement; (iii) formalize the finite Walsh identities, Hermite quadrature and limit arguments in Lean; (iv) determine the next term in \eqref{eq:main}, which may depend on the interior nodes' positions in their parity grids; (v) investigate arithmetic lower bounds rather than only the abstract bias obstruction. None of these is claimed settled here.

\begin{thebibliography}{9}
\bibitem{vendor} ProveIt Contributors, \emph{Vendored grh-lean, readable edition}. \href{https://github.com/VladimirReshetnikov/ProveIt/tree/main/Analysis/GeneralizedRiemannHypothesis/vendor/grh-lean}{Repository directory}, inspected 27 November 2026. Upstream commit as pinned above.
\bibitem{upstream} L. Sernau and collaborators, \emph{The generalized Riemann hypothesis for Dirichlet L-functions, in Lean 4}, rewriting the Alp\"oge proof developed with Claude. \href{https://github.com/lsernau/grh-lean/tree/72fead102f663d5c7a407964a13dc0cf9d977207}{Pinned upstream repository}.
\bibitem{article} ProveIt Contributors, \emph{Beyond the Riemann hypothesis for Dirichlet L-functions: consequences, transfers and limits}, 26 November 2026, \href{https://github.com/VladimirReshetnikov/ProveIt/tree/main/Analysis/GeneralizedRiemannHypothesis/Article}{merged article}; \href{https://github.com/VladimirReshetnikov/ProveIt/blob/main/Analysis/GeneralizedRiemannHypothesis/Article/sections/dir.tex#L3538}{Hermite threshold section}, \href{https://github.com/VladimirReshetnikov/ProveIt/blob/main/Analysis/GeneralizedRiemannHypothesis/Article/sections/dir.tex#L3847}{open questions}. Blob of \texttt{dir.tex} inspected: \texttt{da58d506e24a05e1b888da2802398f277d9ea727}.
\bibitem{bellaiche} J. Bella\"iche, Th\'eor\`eme de Chebotarev et complexit\'e de Littlewood, \emph{Ann. Sci. \`Ec. Norm. Sup\'er.} 49 (2016), 579--632. \href{https://doi.org/10.24033/asens.2291}{doi:10.24033/asens.2291}.
\bibitem{dilv} H. Derksen, P. Ivanov, C. H. Lee, E. Viola, \emph{Pseudorandomness, symmetry, smoothing: II}, \href{https://arxiv.org/abs/2407.12110}{arXiv:2407.12110}, 2024; ECCC TR25-021, 2025.
\bibitem{mv} H. L. Montgomery, R. C. Vaughan, \emph{Multiplicative Number Theory I: Classical Theory}, Cambridge Studies in Advanced Mathematics 97, Cambridge University Press, 2007 (explicit formula and integrated character sums, Ch.~13).
\end{thebibliography}
\end{document}
